\section{Results}\label{sec:results}

\begin{figure}[t]
\centering
\includegraphics[width=\linewidth]{figures/fig1_main_sweep.pdf}
\caption{\textbf{E1.} Honest-receiver accuracy against a coherent lying bloc, as a function of the Byzantine fraction $f$ (TEST split, 5 seeds, full visibility). Majority and label-free Dawid--Skene fall to 0 at $f\ge1/2$ (label switching, \cref{prop:switching}). \aip\ collapses on binary BoolQ at $f=0.5$ and everywhere at $f=0.9$. \race\ never falls below the receiver alone at any $f$, except on MATH-500 at small $f$ (weak-anchor capture, \cref{sec:limits}).}\label{fig:main}
\end{figure}

\subsection{E1: no breakdown point against a coherent bloc}
\Cref{fig:main} is the headline. Averaged over the six benchmarks at $f=0.5$, \race\ reaches \mainRaceFive\% against \mainAipFive\% for \aip, \mainMajFive\% for majority vote, \mainDsFive\% for label-free Dawid--Skene and \mainSelfFive\% for the receiver alone. At $f=0.9$ the figures are \mainRaceNine\%, \mainAipNine\%, \mainMajNine\%, \mainDsNine\% and \mainSelfNine\%. The mean accuracy of majority vote and label-free Dawid--Skene first falls below half the receiver's at $f=\breakdownMaj$ and $f=\breakdownDs$. For \aip\ this happens at $f=\breakdownAip$, and for \race\ at \breakdownRace\ of the tested fractions. The known-channel oracle bounds the remaining headroom at \mainOracleFive\% ($f=0.5$).

The binary case, which \aip\ cannot handle, is where inversion is most valuable. On BoolQ at $f=0.7$, \race\ scores \mainBoolqRaceSeven\% (\aip: \mainBoolqAipSeven\%, receiver: \mainBoolqSelfSeven\%). A bloc that always says the opposite is a perfect oracle once it is read backwards. With all agents honest ($f=0$), \race\ costs nothing: \mainRaceZero\% against \mainAipZero\% for \aip\ and \mainMajZero\% for majority. Across all $6\times10$ cells (Holm-corrected Wilcoxon, \cref{tab:wtl}), \race\ is significantly better than \aip\ in \mainWinsAip\ cells and worse in \mainLossesAip. Against majority it wins \mainWinsMaj\ cells and loses \mainLossesMaj. Against the receiver alone it wins \mainWinsSelf\ and never loses. Every significant loss is on MATH-500 at $f\le0.3$, the weak-anchor regime analysed in \cref{sec:limits}. Clone tempering has a price against a \emph{perfectly} coherent bloc. Counting the bloc's identical vetoes separately happens to help there (\mainRaceNocloneFive\% without tempering at $f=0.5$), but the tempered rule is the one that wins on the confirmation sets (\cref{sec:llm}).

\subsection{E2: the attack zoo and the gate-aware adversary}
\begin{figure}[t]
\centering
\includegraphics[width=\linewidth]{figures/fig2_attack_zoo.pdf}
\caption{\textbf{E2.} Accuracy minus the receiver alone (percentage points), averaged over six benchmarks and three seeds, for 14 attack settings at three Byzantine fractions. Blue: pooling helps. Red: pooling hurts.}\label{fig:zoo}
\end{figure}
\Cref{fig:zoo} organises the attacks by the assumption they target. Three findings stand out.

\emph{Coherence gates are blind to incoherent liars.} Against the gate-aware attacker at $f=0.7$ (coordination $p\le0.5$, i.e.\ below \aip's published ceiling), \aip\ scores \zooGateAipSeven\% and its soft gate \zooGateSoftSeven\%, while \race\ scores \zooGateRaceSeven\% (\cref{fig:gate}). Independent liars, who never coordinate at all, give \zooIndepAipSeven\% versus \zooIndepRaceSeven\%. This is \cref{prop:quantisation} in data. The attacker's coherence is tunable, but its dependence on the truth is not.

\emph{The minimax attack costs little.} Uninformative liars are the minimax strategy of \cref{thm:silence}. Against them \race\ stays at \zooUninfRaceSeven\% at $f=0.7$, close to the honest-only oracle (\zooUninfOracleSeven\%) and above the receiver alone (\zooUninfSelfSeven\%).

\emph{\race's own failure modes.} Camouflage (\cref{sec:limits}) is the one stationary attack that pushes \race\ below the receiver at $f=0.7$ (\zooCamoRaceSeven\% vs \zooCamoSelfSeven\%). The sleeper defeats every static fit.

\begin{figure}[t]
\centering
\includegraphics[width=\linewidth]{figures/fig3_gate_aware.pdf}
\caption{The gate-aware adversary. The upper axis shows the attacker's pair coherence $q(p)$. \aip\ is defeated exactly in the sub-ceiling band. \race\ is nearly flat in $p$ (within about 3.5 points at $f=0.7$), because the attacker's accuracy, not its coherence, sets the weight.}\label{fig:gate}
\end{figure}

\paragraph{Adaptive attackers.} For every defender, benchmark, $f$ and seed we let the attacker pick, on VALIDATION, whichever stationary zoo attack hurts that defender most, and score TEST under that choice (\cref{fig:br}). Under its own best response, \race\ holds \brRaceFive\% at $f=0.5$ and \brRaceSeven\% at $f=0.7$. The corresponding figures are \brAipFive\% and \brAipSeven\% for \aip, \brMajFive\% and \brMajSeven\% for majority, and \brSelfFive\% and \brSelfSeven\% for the receiver alone. Including the sleeper, every static method fails at $f\ge0.5$. \Cref{sec:online} addresses this with online refitting.

\begin{figure}[t]
\centering
\includegraphics[width=\linewidth]{figures/fig9_best_response.pdf}
\caption{Best-response attacker. The attack is selected on VALIDATION for each defender separately and scored on TEST.}\label{fig:br}
\end{figure}

\subsection{E3: replayed LLM deceivers (confirmation set)}\label{sec:llm}
The four prompt-injected attacks of the original study were written by real LLMs. Their lies are neither perfectly wrong nor perfectly coordinated. This set was not inspected before the method was frozen. \race\ is above every non-\race\ method at every $f$, and within $0.1$ points of its own best variant (\cref{fig:llm}). At $f=0.7$ it scores \llmRaceSeven\%, against \llmAipSeven\% for \aip, \llmMajSeven\% for majority, \llmDsSeven\% for Dawid--Skene and \llmSelfSeven\% for the receiver alone. That is within \llmOracleSeven\% of the known-channel oracle and above the honest-only oracle (\llmHonestOracleSeven\%), because inverting the liars adds information that honest agents alone do not have. Across \llmCells\ Holm-corrected comparisons (5 baselines $\times$ 6 benchmarks $\times$ 4 prompts $\times$ 4 fractions), \race\ has \llmWins\ significant wins and \llmLosses\ losses. The confirmation result proper is that of the frozen v3.0 rule, which differs from v3.1 only on BoolQ: \llmWinsVthree\ wins and \llmLossesVthree\ losses. Both superseded variants of the clone rule score lower on this set (\llmRaceRawSeven\% and \llmRaceNocloneSeven\% at $f=0.7$). This is independent support for the one post-hoc design change (\cref{app:devlog}).

\begin{figure}[t]
\centering
\includegraphics[width=\linewidth]{figures/fig4_llm_liars.pdf}
\caption{\textbf{E3.} Replayed real-LLM deceivers (Llama-3.2-3B and Ministral-3-14B under four deception prompts). Confirmation set, evaluated after the method was frozen.}\label{fig:llm}
\end{figure}

\subsection{E7: does RACE estimate the channels?}
\race's estimates are label-free. We compared them with each peer's true accuracy on HISTORY, which only the evaluator can compute, over \chanN\ receiver--peer channels spanning seven attacks, four fractions and six benchmarks (\cref{fig:chan}). The correlation is $r=\chanCorr$ and the mean absolute error is \chanMae. \race\ inverts \chanByzInvRace\% of Byzantine channels, against \chanByzInvAip\% for \aip. It falsely inverts \chanHonInvRace\% of honest channels (\aip: \chanHonInvAip\%). The Byzantine channels \race\ trusts (\chanByzTrustRace\%) are liars that tell the truth often enough to be positively informative, such as partial and camouflage attackers and LLM liars that fail to lie.

\begin{figure}[t]
\centering
\includegraphics[width=\linewidth]{figures/fig5_channel_estimation.pdf}
\caption{\textbf{E7.} Left: label-free \race\ estimates versus true channel accuracy (evaluator-side). Right: share of TRUST/\allowbreak DISCARD/\allowbreak INVERT decisions for honest and Byzantine peers.}\label{fig:chan}
\end{figure}

\paragraph{Anchor strength.} \Cref{fig:competence} groups receivers by their own accuracy. For receivers above 50\%, \race\ ends more than 5 points worse than trusting themselves in only \riskAboveRace\% of cases. The rates are \riskAboveAip\% for \aip\ and \riskAboveMaj\% for majority. For receivers at or below 50\%, the weak-anchor regime of \cref{prop:switching}, the rate rises to \riskBelowRace\%.

\begin{figure}[t]
\centering
\includegraphics[width=\linewidth]{figures/fig10_receiver_competence.pdf}
\caption{Receiver competence. Left: mean accuracy. Right: share of receivers that end $>5$ points worse than trusting themselves.}\label{fig:competence}
\end{figure}

\subsection{E5: how much history?}
\race\ needs little history (\cref{fig:history}). Averaged over attacks and benchmarks at $f=0.5$, it reaches \histRaceFive\% with 5 history questions and \histRaceTwoZero\% with 20. \aip\ reaches \histAipFive\% and \histAipTwoZero\%: its significance test has a minimum-support floor. This matches \cref{prop:history}. An always-wrong liar is far from chance, so few joint observations separate it.

\begin{figure}[t]
\centering
\includegraphics[width=\linewidth]{figures/fig6_history_length.pdf}
\caption{\textbf{E5.} Accuracy versus history length.}\label{fig:history}
\end{figure}

\subsection{E4: replication, composition and swarm size}
% RESULTS-DEPENDENT: interpretation below must be checked against results/swarm before release.
Deployed swarms usually replicate one model (\cref{fig:swarm}). In a swarm of ten replicas of Llama-3.2-3B, the weakest roster model, facing a coherent bloc at $f=0.5$, \race\ scores \swarmHomLlamaRace\% (\aip: \swarmHomLlamaAip\%, receiver alone: \swarmHomLlamaSelf\%). With ten replicas of Qwen-3.8-27B the figures are \swarmHomQwenRace\%, \swarmHomQwenAip\% and \swarmHomQwenSelf\%. In a homogeneous swarm the honest agents are one source. Clone tempering counts them once, and the only information left to pool is what the liars leak. Adding the two weak-tier models to the roster (``7 + weak arm'') gives \swarmWeakRace\% for \race\ against \swarmWeakAip\% for \aip. Scaling the frozen swarm from 5 to 40 agents at fixed $f$ changes \race\ from \sizeRaceFive\% to \sizeRaceFourZero\% and \aip\ from \sizeAipFive\% to \sizeAipFourZero\% (pooled over $f\in\{.3,.5,.7\}$ and the coherent and gate-aware attacks). Both results are consistent with \cref{thm:silence}: extra copies of the same information do not add information.

\begin{figure}[t]
\centering
\includegraphics[width=\linewidth]{figures/fig7_swarm.pdf}
\caption{\textbf{E4.} Left: composition and replication at $f=0.5$ (coherent bloc). Right: swarm size.}\label{fig:swarm}
\end{figure}


\subsection{E9: what are the liars worth?}\label{sec:budget}
\Cref{thm:silence} says liars cannot remove information from a receiver that knows the joint law. It does not say how much information they add, or how much of it a label-free receiver can recover. We measure both with two oracles that see the gold labels of the HISTORY split. The \emph{liar-removal oracle} knows who lies, removes every liar, and decodes the honest agents with their true accuracies and true source groups. The \emph{known-channel oracle} keeps every agent, with true accuracies and true source groups. We call the gap between them the \emph{information budget} of an attack. E9 evaluates both oracles, \race\ and the extension below on \nWorldsExt\ worlds: 14 symbolic and 4 LLM attack settings, $f\in\{0.3,0.5,0.7\}$, six benchmarks and three seeds. \Cref{fig:budget} shows every method relative to the liar-removal oracle.

\paragraph{Truth-dependent liars are worth more than their absence.} For independent and gate-aware liars the budget is positive and grows with $f$. For independent liars it is \budgetIndepThree, \budgetIndepFive\ and \budgetIndepSeven\ points at $f=0.3$, $0.5$ and $0.7$. At $f=0.7$, three honest agents plus seven liars decode better than the three honest agents alone. \race\ captures most of this budget without labels and without knowing who lies. At $f=0.7$ it scores \extRaceIndepSeven\% against independent liars, while the oracle that removes every liar scores \extRemovalIndepSeven\%. Pooling tasks over benchmarks, \race\ is significantly \emph{above} the liar-removal oracle in \pooledRemovalWins\ of \pooledRemovalCells\ (attack, $f$) cells. All of these cells are truth-dependent attacks: independent, gate-aware with $p<1$ at every $f$, and the coherent bloc and attractor at $f=0.7$. A rule that only filters suspected liars sees at most the honest reports, so in expectation it cannot beat the Bayes receiver on those reports, for which this oracle is a proxy. Exceeding it requires reading the liars.

\paragraph{LLM deceivers: the value is neutralisation, not extra information.} For the four replayed LLM attacks the budget is close to zero: \budgetLlmThree, \budgetLlmFive\ and \budgetLlmSeven\ points on average. Deception prompts produce lies that are only weakly tied to the truth. At the one-coin level they carry little beyond what the honest agents already say. \race's contribution is then to cancel the damage. At $f=0.7$, majority vote falls to \extLlmMajSeven\% while \race\ scores \extLlmRaceSeven\%. That recovers \recoveryLlmSeven\% of the gap between the receiver alone (\extLlmSelfSeven\%) and the liar-removal oracle (\removalSeven\%); at $f=0.5$ it recovers \recoveryLlmFive\%. In the per-benchmark cells, \race\ is statistically indistinguishable from the liar-removal oracle in \removalLlmTies\ of \removalLlmCells\ LLM cells and below it in \removalLlmLosses.

\paragraph{A negative budget certifies a model violation.} The known-channel oracle assumes conditionally independent, symmetric channels. An attack against which it does \emph{worse} than the liar-removal oracle therefore lies outside that model. Three attacks do: echo (\budgetEchoSeven\ at $f=0.7$), camouflage (\budgetCamoSeven) and the sleeper (\budgetSleeperFive\ at $f=0.5$). The echo attacker copies honest reports, so its errors are not independent given the truth. The camouflage attacker conditions on the honest vote. The sleeper is non-stationary. These are exactly \race's failure modes (\cref{sec:limits}). An evaluator with labels can therefore use the sign of the budget to check whether a new attack is inside the model class that \race\ assumes. Uninformative liars have a budget of essentially zero (\budgetUninfSeven\ at $f=0.7$), as \cref{thm:silence} predicts for the minimax attack. \race\ pays a small estimation cost there (\overRemovalUninfSeven\ points relative to the liar-removal oracle), which is the price of not knowing who is silent.

\paragraph{A negative result: conditioning on disagreement (\raceD).} Camouflage liars earn trust on questions where every honest agent agrees. A natural fix is to fit the channels only on history questions where the heard answers disagree, falling back to all questions when fewer than five qualify. We call this \raceD. On the 14 attack settings other than camouflage, echo and sleeper, \raceD\ never differs significantly from \race\ (\raceDOtherNonTies\ non-ties in \raceDOtherCells\ cells). Against camouflage it helps at $f=0.7$ (\raceDCamoSeven\ points on average) but hurts at $f=0.3$ (\raceDCamoThree). Overall it wins \raceDCamoWins\ and loses \raceDCamoLosses\ of \raceDCamoCells\ camouflage cells. It also hurts against echo (\raceDEchoSeven\ at $f=0.7$), and against sleepers at $f=0.3$ (\raceDSleeperThree). Across all \raceDCells\ cells it wins \raceDWins\ and loses \raceDLosses. Restricting to disagreement removes the questions that pin down the honest peers' accuracies, trading one misspecification for another. We keep the frozen \race\ as the method. \raceD\ was specified and implemented before E9 was run, and \race\ was not modified afterwards (\cref{app:devlog}).

\begin{figure}[t]
\centering
\includegraphics[width=\linewidth]{figures/fig12_information_budget.pdf}
\caption{\textbf{E9.} What the liars are worth. Accuracy minus the liar-removal oracle, which knows who lies, removes them, and decodes the honest agents with their true accuracies. Violet stars: the known-channel oracle, which keeps the liars; its offset is the information budget. Filled blue: \race. Hollow blue: \raceD. Shaded: better than removing every liar. Means over six benchmarks and three seeds (TEST split). The sleeper is omitted because a history-fitted oracle is invalid by construction; see \cref{tab:budget}.}\label{fig:budget}
\end{figure}


\subsection{E11: classical crowdsourcing estimators}\label{sec:crowd}
Is the receiver anchor necessary, or would any mature crowdsourcing estimator do? E11 adds four classical label-free estimators to both Dawid--Skene variants: iterative weighted majority voting (IWMV)~\citep{li2014iwmv}, MACE~\citep{hovy2013}, GLAD~\citep{whitehill2009} and, on binary questions, the message-passing estimator of Karger, Oh and Shah (KOS)~\citep{karger2014}. The protocol is exactly that of \race. Each honest agent runs every estimator on its own HISTORY window of peer answers, no estimator sees a label, and ties go to the agent's own answer. E11 covers \nWorldsCrowd\ worlds: six benchmarks; a coherent bloc at $f\in\{0.1,0.3,0.5,0.7\}$; gate-aware liars ($p=0.25$) and the four LLM prompts at $f\in\{0.3,0.5,0.7\}$; and three seeds. \Cref{fig:crowd} shows the result.

\paragraph{With a liar minority the classical estimators are as good as \race.} At $f=0.3$ the best of the four estimators defined on every benchmark (IWMV, MACE, GLAD, label-free Dawid--Skene) scores \crowdBestClassicalThree\%, and \race\ \crowdRaceThree\%. Over the \crowdClassicalCellsLow\ comparisons between \race\ and a classical estimator at $f\le0.3$, \race\ ties in \crowdClassicalTiesLow\ and wins \crowdClassicalWinsLow. This is the regime these estimators were designed for, and they work.

\paragraph{With a liar majority every one of them collapses.} At $f=0.7$ the best of those four scores \crowdBestClassicalSeven\%, full-confusion Dawid--Skene \crowdDsFullSeven\% on the four categorical benchmarks, and majority vote \crowdMajSeven\%. \race\ scores \crowdRaceSeven\%, the known-channel oracle \crowdOracleSeven\% and the receiver alone \crowdSelfSeven\%. On BoolQ, KOS scores \crowdBoolqKosSeven\% against \race's \crowdBoolqRaceSeven\%. The failure is the label-switching symmetry of \cref{prop:switching}, not a weakness of any single estimator. The likelihood of a latent-class model is invariant to relabelling the truth. Without an anchor, EM settles on the relabelling in which the \emph{majority} is reliable, so it decodes the liars' labelling faithfully and inverts every honest agent. The receiver anchor breaks exactly this symmetry. Anchoring Dawid--Skene and nothing else gives \race\ without clone tempering, which E1--E3 already evaluate.

\paragraph{Paired tests.} There are \crowdClassicalCells\ (benchmark, $f$, attack) comparisons between \race\ and a classical estimator. \race\ wins \crowdClassicalWins\ and loses \crowdClassicalLosses, with Holm correction over all \crowdCells\ comparisons of the study. At $f\ge0.5$ it wins \crowdClassicalWinsHigh\ of \crowdClassicalCellsHigh. The classical estimators are therefore not stronger baselines than those of E1--E3. Their failure at $f\ge 1/2$ is the reason an anchored estimator is needed.

\begin{figure}[t]
\centering
\includegraphics[width=\linewidth]{figures/fig14_crowd_baselines.pdf}
\caption{\textbf{E11.} Classical crowdsourcing estimators against \race. (a)~Accuracy of an honest agent as the liar share $f$ grows. The yellow band spans IWMV, MACE, GLAD and label-free Dawid--Skene; the dashed line is the best of them. (b)~Accuracy per benchmark at $f=0.7$. KOS is defined for binary questions only, and full-confusion Dawid--Skene is not run on open-answer benchmarks. Means over attacks and seeds, TEST split.}\label{fig:crowd}
\end{figure}


\subsection{E6: changing attackers}\label{sec:online}
The sleeper is the one attack in the zoo that defeats every static defender, including the known-channel oracle, because it violates stationarity. We therefore re-ran five attacker schedules under a strictly causal protocol on MMLU, MedQA, BoolQ and MATH-500 ($f\in\{0.3,0.5,0.7\}$, two seeds; 120 streams of 200 questions). At each step every honest receiver predicts from a state fitted only on earlier questions, then commits the observation. Methods refit every 8 questions after a 40-question warm-up. A test confirms that future broadcasts or labels cannot change a past prediction.

\paragraph{Sleepers.} The sleeper answers honestly for 100 questions and then lies coherently. We average over the last 60 questions, which start 40 questions after the switch, and over the three values of $f$. A \race\ fit frozen at warm-up scores \onSleeperFixed\%. Cumulative refitting recovers to \onSleeperCum\%. Forgetting ($\gamma=0.97$) recovers to \onSleeperDecay\% and a 64-question window to \onSleeperWin\%, both above the receiver alone (\onSleeperSelf\%) and far above a windowed \aip\ (\onSleeperAipWin\%) and majority vote (\onSleeperMaj\%). At $f=0.7$ the frozen fit and majority vote both score \onSleeperSevenFixed\%, while forgetting scores \onSleeperSevenDecay\% (receiver alone: \onSleeperSevenSelf\%). Once the sleeper starts lying, it is a coherent bloc like any other, and the forgetting receiver learns to read it backwards within about 40 questions (\cref{fig:online}).

\paragraph{Regime switches and the price of forgetting.} When the bloc switches from coherent to independent lies, or back, forgetting scores \onCohIndDecay\% and \onIndCohDecay\% after the switch (cumulative: \onCohIndCum\% and \onIndCohCum\%; windowed \aip: \onCohIndAipWin\% and \onIndCohAipWin\%). A bloc that alternates between coherent lies and silence every 40 questions leaves \race\ at \onToggleDecay\%, against \onToggleMaj\% for majority vote. On a stationary coherent schedule, forgetting costs nothing measurable (\onStatDecay\% vs.\ \onStatCum\% cumulative). Forgetting is therefore the default we recommend when stationarity cannot be assumed, and it is exposed as \texttt{TrustLayer(forgetting=0.97)}.

\begin{figure}[t]
\centering
\includegraphics[width=\linewidth]{figures/fig8_online.pdf}
\caption{\textbf{E6.} Causal online accuracy (16-question rolling mean) at $f=0.5$ under four non-stationary attacker schedules, averaged over four benchmarks and two seeds. Dashed lines: the attacker's switches. Shaded (right): the bloc is silent.}\label{fig:online}
\end{figure}


\begin{table}[t]
\centering\small
\caption{Significant wins/ties/losses of \race\ against each baseline (Holm-corrected Wilcoxon signed-rank and paired bootstrap interval must agree, $\alpha=0.05$, task as unit).}\label{tab:wtl}
\begin{tabular}{llrrrr}
\toprule
 & verdict & Wins & Ties & Losses & Cells \\
Study & Baseline &  &  &  &  \\
\midrule
\multirow[t]{5}{*}{E3 LLM deceivers} & AIP gated & 37 & 59 & 0 & 96 \\
 & AIP soft gate (M1) & 28 & 68 & 0 & 96 \\
 & Dawid–Skene (label-free) & 35 & 61 & 0 & 96 \\
 & Majority vote & 37 & 59 & 0 & 96 \\
 & Receiver alone & 59 & 37 & 0 & 96 \\
\cline{1-6}
\multirow[t]{5}{*}{E1 coherent sweep} & AIP gated & 20 & 38 & 2 & 60 \\
 & AIP soft gate (M1) & 14 & 44 & 2 & 60 \\
 & Dawid–Skene (label-free) & 31 & 26 & 3 & 60 \\
 & Majority vote & 41 & 18 & 1 & 60 \\
 & Receiver alone & 38 & 22 & 0 & 60 \\
\cline{1-6}
\multirow[t]{5}{*}{E2 attack zoo} & AIP gated & 100 & 147 & 5 & 252 \\
 & AIP soft gate (M1) & 73 & 175 & 4 & 252 \\
 & Dawid–Skene (label-free) & 87 & 157 & 8 & 252 \\
 & Majority vote & 122 & 127 & 3 & 252 \\
 & Receiver alone & 136 & 96 & 20 & 252 \\
\cline{1-6}
\bottomrule
\end{tabular}

\end{table}
